\documentclass[12pt]{article}
\usepackage[utf8]{inputenc}
\usepackage[T1]{fontenc}
\usepackage[italian]{babel}
\usepackage{multicol}
\usepackage{enumitem}
\usepackage{array}
\usepackage{graphicx}
\usepackage{tikz}
\usetikzlibrary{
    shapes.geometric,
    arrows.meta,
    positioning,
    calc,
    patterns,
    decorations.pathmorphing
}
\usepackage[table]{xcolor}
\usepackage[margin=2cm]{geometry}
\usepackage{amsmath,amssymb}
\setlength{\parindent}{0pt}
\definecolor{headercolor}{RGB}{220,220,220}
\newcounter{quesito}
\setcounter{quesito}{0}
\newcommand{\nuovoquesito}{\stepcounter{quesito}\textbf{Domanda \arabic{quesito}.}}
\newcommand{\itemdomanda}[2]{%
    \par\noindent
    \begin{minipage}[t]{\linewidth}
        \nuovoquesito \\ #1
        \begin{enumerate}[label=(\Alph*), nosep, leftmargin=*]
            #2
        \end{enumerate}
    \end{minipage}
    \vspace{10pt}
}
\begin{document}
\section*{Simulazione}
\hfill Data: 16 apr 2026

\itemdomanda{Se A è più alto di B e B è più alto di C, allora:}{
    \item C è più alto di A
    \item A è più basso di C
    \item A è più alto di C
    \item B è il più alto di tutti
    \item Non si può stabilire nulla
}

\itemdomanda{Date le seguenti premesse:
    \begin{enumerate}
        \item Nessun avaro è altruista.
        \item Solo gli avari tengono da conto i gusci d'uovo.
    \end{enumerate}
    Quale conclusione se ne può trarre?}{
    \item Nessun altruista tiene da conto i gusci d'uovo
    \item Chi tiene da conto i gusci d'uovo è altruista
    \item Tutti gli altruisti sono avari
    \item Qualche avaro è altruista
    \item Gli altruisti tengono da conto solo i gusci d'uovo
}

\itemdomanda{Dire quante stelle sono comprese sia nel triangolo, sia nel cerchio ma non nel quadrato e neppure nel rettangolo.
\begin{center}
\begin{tikzpicture}[scale=0.35]
    % Forme geometriche (approssimate dal disegno originale)
    \draw[thick] (0,5) circle (4); % Cerchio
    \draw[thick] (-1,8) -- (12,12) -- (12,1) -- cycle; % Triangolo grande
    \draw[thick] (2,3) rectangle (10,11); % Rettangolo orizzontale
    \draw[thick] (4,0) rectangle (7,7); % Quadrato/Rettangolo verticale

    % Stelle (posizionate strategicamente per il conteggio richiesto)
    % Area target: Intersezione Cerchio e Triangolo, fuori dagli altri due
    \node at (-0.5,6) {$\star$};
    \node at (0.5,4.5) {$\star$};
    \node at (-1,4) {$\star$};
    
    % Altre stelle sparse per riempire il disegno come l'originale
    \node at (-2,5) {$\star$}; \node at (5,9) {$\star$}; \node at (8,4) {$\star$};
    \node at (10,10) {$\star$}; \node at (1,1) {$\star$}; \node at (11,6) {$\star$};
    \node at (3,6) {$\star$}; \node at (6,2) {$\star$}; \node at (9,8) {$\star$};
    \node at (4,10) {$\star$}; \node at (1,8) {$\star$}; \node at (5,5) {$\star$};
\end{tikzpicture}
\end{center}
}{
    \item 2
    \item 3
    \item 4
    \item 5
    \item 6
}

\itemdomanda{$ \sqrt[4]{16} $ è uguale a:
}{
\item $ 4 $
\item $ 8 $
\item $ 2 $
\item $ 1 $
\item $ \frac{1}{4} $
}

\itemdomanda{$ \sqrt{x} > \sqrt{y} $ è verificata
}{
\item per $ x > y $
\item per $ x > y $ solo se $ x > 1 $ e $ y \ge 1 $
\item per $ x > y $ solo se $ x > 0 $ e $ y \ge 0 $
\item mai
\item sempre
}

\itemdomanda{La disequazione $ \frac{1}{x} > 0 $ è verificata per:
}{
\item $ x \neq 0 $
\item $ x < 0 $
\item $ x > 0 $
\item $ \forall x \in \mathbb{R} $
\item $ x \geq 1 $
}

\itemdomanda{L'equazione $ x^4-5x^2+4=0 $ ha soluzioni reali:
}{
\item $ x=\pm 1 $
\item $ x=\pm 2 $
\item $ x=\pm 1, x=\pm 2 $
\item $ x=1, x=4 $
\item non ha soluzioni reali
}

\itemdomanda{Nel campo dei numeri reali, il numero degli zeri del polinomio $P(x) = x^{2} + a^{2}$ con $a \in \mathbb{R}$ è
}{
\item dipende da $a$ se $a \neq 0 $
\item pari a $2$ se $a = -1$
\item pari a $2$ se $a = 1$
\item pari a $2$ se $a$ è un quadrato perfetto
\item sempre pari a zero se $a \neq 0$
}

\itemdomanda{Velocità luce nel vetro $1.97 \times 10^8$ m/s. Indice rifrazione?}{
    \item 0.88
    \item 1.34
    \item 1.52
    \item 1.66
    \item 1.88
}

\itemdomanda{Fenomeno ondulatorio NON evidenziabile con onde sonore:}{
    \item Rifrazione
    \item Interferenza
    \item Polarizzazione
    \item Riflessione
    \item Diffrazione
}

\end{document}
